\fancyhead[L]{IHM11}
\renewcommand{\arraystretch}{1}


\titre{}{IHM11 - HTML et CSS}


 Vers la fin des années 1950, face aux difficultés à pouvoir échanger des informations rapidement et à distance, quelques chercheurs mettent au point un outil capable de faire communiquer deux ordinateurs entre eux (client/serveur), reliés par une ligne téléphonique : c'est la naissance du premier réseau internet. 


\section{Le langage HTML}

\subsubsection{Historique}

\begin{minipage}{12cm}
 Pour partager des informations avec ses collègues chercheurs au CERN, Tim Berners-Lee met au point dans les années 1990 un système de liens (ou hyperliens) reliant les informations contenues dans les ordinateurs entre elles grâce à des adresses web ou \textbf{URL (Uniform Ressource Locator)} : c'est la naissance du \textbf{Web (world wide web)}.
 
\end{minipage}
\begin{minipage}{6cm}
\begin{center}
\includegraphics[scale=0.5]{client-serveur.ps}
\end{center}
\end{minipage}

\medskip

  Le \textbf{world wide web} (abrégé en « www » ou « le Web ») est le réseau mondial (ou « la Toile ») constituée de l'ensemble des pages web liées entre elles par des liens hypertextes. Afin de structurer les données présentes sur une page web, Tim Berners-Lee développa un langage de description dans lequel sont écrite les pages web: le HTML (hypertext markup language). 


 Tout document texte ayant l'extension \verb+.htm+ ou \verb+.html+, constitue une page web. Cependant, pour que cette page soit lisible par un navigateur (Chrome, Edge, Firefox, Opera ou Safari), elle doit vérifier les règles, le \textbf{balisage du langage HTML}. 
 
 \subsubsection{Utilisation}
 
  Les informations d'une page HTML sont encadrées par différents balises : une ouvrante \verb+<html>+ et une fermante \verb+</html>+. Toute balise qui a été ouverte doit être fermée sauf les balises orphelines. Une page HTML 5 est construite sur la base suivante : 
  
  
  \begin{center}
  \includegraphics[scale=0.4]{ScreenShot_20240704152929.ps}
  \end{center}
  
  
  Toutes les balises HTML et leurs attributs sont présentés sur \href{https://www.w3schools.com/tags/ref_byfunc.asp}{w3school}
  
  
  \newpage
  
  
  \exer Structure de base en HTML 5
  
  \begin{enumerate}
  	\item Ouvrir un éditeur de texte (\verb%Notepad++ % ou autre) et sélectionner, pour \verb%Notepad++%, dans l'onglet Langage/HTML et Encodage/Encoder en \verb%UTF-8%.
  	
\item Recopier la structure de base d'une page HTML donné ci-dessus en personnalisant le titre et le contenu. Sauvegarder sous le nom \og page1.html \fg et ouvrir votre fichier avec le navigateur Firefox. Où s'affiche votre titre ? Où s'affiche votre contenu ? \dotfill

\dotfill

\item Depuis le navigateur, ouvrir / Plus d'outils / Outils de développement (Ctrl + Maj + I) pour sélectionner l'onglet "Éléments" pour voir le code source de votre page. Modifier son contenu depuis le navigateur. 

\item Que se passe-t-il si vous modifier le code source dans le navigateur ? Votre fichier d'origine (sur le serveur) est-il modifié ? \dotfill

\dotfill
  \end{enumerate}
  
  \medskip
  
  \cadreR{Pour la suite, aidez-vous du mémo HTML}
  
  \bigskip
  
  \exer Structuration du contenu en HTML
  
  À chacune des étapes, sauvegarder votre fichier et actualiser son affichage dans le navigateur pour voir le résultat.
  
  \begin{enumerate}
  	\item Ajouter à votre page un titre de niveau 1 en utilisant la balise \verb%<h1>%. N'oublier pas de refermer votre titre avec \verb%</h1>%.

Par défaut, comment apparaissent les titres de niveau 1 ? \dotfill
	\item Ajouter un paragraphe \verb%<p>% avec un long texte (copier-coller ce que vous voulez). N'oublier pas de refermer votre balise.
	 
	 Par défaut, les paragraphes s'adaptent-ils à la largeur de la page ? \dotfill
	\item Ajouter un deuxième titre de niveau 1 puis un titre de niveau 2 \verb%<h2>%.
	\item Ajouter une liste non-ordonnée \verb%<ul>% qui contiendra trois éléments \verb%<li>%.
	
Par défaut, quelle puce apparaît devant chaque élément de la liste ? \dotfill
	\item Convertir votre liste non-ordonnée en liste ordonnée avec la balise \verb%<ol>%.
	
Par défaut, comment s'insère tous ses éléments dans votre page ? côte à côte ou les uns sous les autres ? 

\dotfill
  \end{enumerate}
  
  
  \medskip
  
  \exer Les liens hypertextes
  
  \begin{enumerate}
  	\item     Ajouter un lien vers \verb%https://www.w3schools.com% avec la balise \verb%<a>% à l'intérieur d'un des paragraphes. Il faut pour cela renseigner l'attribut \verb%href% à l'intérieur de la balise (voir mémo) avec comme valeur, l'URL du lien.
  	
    Par défaut, comment s'insère le lien dans le paragraphe ? au sein du texte ou en-dessous ? \dotfill
    
    \dotfill
    
    \item Ajouter une info-bulle en renseignant l'attribut title dans la balise ouvrante de votre lien. Qu'observe-t-on en passant la souris au-dessus du lien ?
    
    \dotfill


  \end{enumerate}
  
  \medskip
  
  \exer Les images
  
  \begin{enumerate}
  	\item     Trouver une image libre de droit en ligne et copier son URL (CTRL+C).
  	
    Sous votre premier titre, insérer cette une image avec la balise orpheline <img>. Renseigner ses attributs \verb+src+ (URL), \verb+alt+ (texte alternatif descriptif de l'image en cas de lien mort).
    
    Par défaut, une image est-elle de type « block » ou « inline » ? \dotfill
    \item Insérer une erreur dans l'URL de votre image. Qu'est-ce qui s'affiche à la place de l'image ?
    
    \dotfill
    \item Ajouter un attribut pour que l'image fasse 200 pixels de large.


  \end{enumerate}
  
  \newpage
  
  \section{Le langage CSS}
  
  
  \subsubsection{Historique}
  
   Depuis 1997 et l'apparition du HTML 4, le choix a été fait de séparer le \textbf{fond} (c'est-à-dire le contenu) et la \textbf{forme} (c'est-à-dire l'aspect) d'une page web. Cela permet d'apporter des modifications à l'aspect visuel d'une page web sans toucher à son contenu.
   
    Ainsi, on utilise en réalité deux langages pour pouvoir créer une page web : 
    
    \begin{liste}
    	\item le langage \textbf{HTML} qui contient le \textbf{fond} (texte, images, videos, liens, ...)
    	\item le langage \textbf{CSS} qui décrit la \textbf{forme} (couleurs, cadres, positions, tailles, police de caractères, ...)
    \end{liste}
    
    

\medskip

\begin{minipage}{14cm}
Le \textbf{CSS (Cascading Style Sheets)} est venu remplacer les anciennes balises de mise en forme qui existait en HTML 1, 2 et 3, car les fichiers HTML devenaient trop lourds et confus. Cela imposait par exemple de modifier une à une toutes les pages d'une site web pour en changer son aspect (sa charte graphique). Désormais, il suffit de modifier le fichier CSS du site web pour que toutes les pages soit mises en forme de la même manière.

\end{minipage}
\begin{minipage}{4cm}
\begin{center}
\includegraphics[scale=0.6]{fichiers.ps}
\end{center}

\end{minipage}



\subsubsection{Utilisation}

\cadreR{Pour la suite, aidez-vous du mémo CSS.}

\bigskip

 Pour lier le fichier HTML à la feuille de style CSS, il doit contenir dans son en-tête \verb%<head>%	, une référence au fichier CSS : 
 
 \begin{lstlisting}[language=HTML,numbers=left]
<link href="feuille1.css" rel="stylesheet">
 \end{lstlisting}
   
   \bigskip
   
   \exer Association d'une feuille de style CSS
   
   \begin{enumerate}
   	\item     Créer un nouveau fichier texte nommé « style1.css » dans le même dossier que votre fichier HTML et l'ouvrir avec un éditeur texte.
    \item Lier votre fichier HTML à ce fichier CSS en ajoutant la bonne balise à l'intérieur de l'en-tête \verb+<head>+ (voir mémo CSS).


   \end{enumerate}
   
   \bigskip
   
   \exer Les sélecteurs
   
   \begin{enumerate}
   	\item À l'aide du sélecteur \verb+h1+, mettre les titres de niveau 1 en gras en copiant dans votre fichier CSS le code : 
   	
   	\begin{lstlisting}[language=HTML,numbers=left]
h1 {
  font-weight: bold;
}
   	\end{lstlisting}
   	
   	\item À l'aide du bon sélecteur, faire que les paragraphes soient de couleur \verb%#0E3C72% (héxadécimal).
Quel sélecteur avez-vous utilisé ?
Combien de couleur peuvent être codée avec six caractères en hexadécimal ? Cela correspond à combien d'octet ? De bits ?
Comment modifier ce sélecteur pour que le style des paragraphes \verb%p% s'applique aussi aux listes (\verb%ul% et \verb%ol%) ? \dotfill

\dotfill

\dotfill
	\item Un de vos paragraphes est une définition et vous souhaitez le faire apparaitre différemment à l'écran. Pour cela ajouter l'attribut \verb%class="definition"% dans la balise ouvrante de ce paragraphe. À l'aide du sélecteur CSS \verb%.definition% le mettre en italique.
Sur ce paragraphe, deux styles s'appliquent : celui défini par le sélecteur \verb%p% et celui du sélecteur \verb%.definition%.
Lequel est prioritaire ? \dotfill
	\item Pour identifier de façon unique une balise, utiliser l'attribut \verb%id="texte"%. Le sélecteur correspondant en CSS est \verb%#texte%.
Appliquer cette méthode de sélection à une de vos balises.
   \end{enumerate}
   
   
   \bigskip
   
    \textbf{Une balise (balise fille) contenue dans une autre (balise mère) hérite de ses propriétés}. C'est le principe en cascade du code CSS.
    
Les \textbf{sélecteurs CSS} permettent de choisir sur quelle balise appliquer un style. Il est possible de pointer les balises par leur nom mais aussi en les regroupant avec un même \textbf{attribut class} ou encore de les pointer de façon unique avec l'\textbf{attribut id.}

Plus le pointeur est précis/sélectif, plus il est prioritaire.\\

Il est toujours possible d'inclure du code CSS sur du HTML en utilisant dans une balise HTML l'attribut \verb+style+.
Cette pratique doit être limitée ou réservée à des mises en forme pérennes. Il est préférable d'utiliser un fichier CSS séparé. 


\begin{lstlisting}[language=HTML,numbers=left]
<p style="color : orange ;">Ce paragraphe apparaît en orange</p>
\end{lstlisting}

Tous les sélecteurs CSS sont présentées ici : \href{https://www.w3schools.com/cssref/css_selectors.php}{w3schools.com : sélecteurs}.

Tous les attributs sont présentés là : \href{https://www.w3schools.com/cssref/index.php}{w3schools : attributs. }


\newpage

\includegraphics[scale=0.25]{ScreenShot_20240704171446.ps}

\includegraphics[scale=0.25]{ScreenShot_20240704171505.ps}

\includegraphics[scale=0.3]{ScreenShot_20240704171532.ps}

\includegraphics[scale=0.3]{ScreenShot_20240704171608.ps}

\includegraphics[scale=0.3]{ScreenShot_20240704171624.ps}

\includegraphics[scale=0.3]{ScreenShot_20240704171633.ps}

%%%%%%%%%%%%%%%%%%%%%%ùùù

\includegraphics[scale=0.66]{ScreenShot_20240704172220.ps}

\includegraphics[scale=0.66]{ScreenShot_20240704172233.ps}

\includegraphics[scale=0.66]{ScreenShot_20240704172244.ps}

\includegraphics[scale=0.66]{ScreenShot_20240704172255.ps}



\import{../../../../preambule}{raz_compteurs.tex}
\newpage